
\textbf{Spurious Correlation Benchmarks.} Waterbirds correlates bird species with background (waterbirds on water 90\% of training samples). Standard ERM achieves 95\% average accuracy but 60\% worst-group accuracy. CelebA exhibits gender-hair color correlation (blonde hair predominantly in female training images). These benchmarks demonstrate spurious features harm minority group performance when the spurious feature is known. Existing detection methods require knowing which feature to search for.

\textbf{Robust Learning Methods.} GroupDRO minimizes worst-group loss via group annotations, achieving 80-85\% worst-group accuracy on Waterbirds (294 GitHub stars). JTT trains a two-stage model: identify hard samples from initial ERM, then upsample and retrain (78\% worst-group accuracy, 72 GitHub stars). SCER refines prototypes in embedding space, reporting estimated 90\% worst-group accuracy (code unavailable for reproduction as of August 2026). GroupDRO and JTT require annotations or two-stage training. SCER operates on embeddings post-hoc, missing gradient-level training dynamics.

\textbf{Gradient Attribution for Spurious Detection.} Grad-CAM produces saliency maps via weighted activation gradients (10,000+ citations). Integrated Gradients computes attribution via path integral from baseline to input. User studies demonstrate practitioners cannot detect unknown spurious correlations using Integrated Gradients or Grad-CAM results alone (109 citations). Key insight: attribution explains which features matter (result interpretation), not whether the process is abnormal (mechanism detection).

GAIA (Gradient Abnormality for Out-of-Distribution Detection) detects distribution shift via gradient zero-deflation and channel variance, achieving FPR95 reduction of 45.41\% on CIFAR100 (16 citations). Operates on gradient process rather than result. SPROD detects spurious-correlated OOD via prototype-based divergence (4 citations). GAIA applied to OOD only, not subpopulation shift. SPROD post-hoc detection only.

We extend GAIA to subpopulation shift, validate spurious-conflict mechanism via synthetic causality tests, and propose spatial gradient regularization as annotation-free mitigation, complementing embedding-based methods by operating on gradient space.
